\subsection{Hilbert-Schmidt 映射的对角化}

定理 2. --- 设 E 与 F 为两个希尔伯特空间，u 为从 E 到 F 中的 Hilbert-Schmidt 映射。存在 E 的标准正交基 \((e_{i})_{i\in I}\)，它被 u 变为 F 中的一个正交族。

用 B 表示 E 的（闭）单位球，并赋以弱化拓扑；这是一个紧空间（V, p.~17）。我们令 \(\mathbf{Q}(x) = \|u(x)\|^{2}\)（对所有 \(x \in B\)）。最后用 P 表示 E 中满足下列性质的向量 x 的全体；

\begin{enumerate}
\def\labelenumi{(\Alph{enumi})}
\setcounter{enumi}{7}
\tightlist
\item
  对每个 \(y \in \mathbf{E}\)（与 \(x\) 正交），元素 \(u(y)\)（属于 \(\mathbf{E}\)）与 \(u(x)\) 正交。
\end{enumerate}

引理 4. --- 函数 \(\mathbf{Q}:\mathbf{B}\to \mathbf{R}\) 是连续的。

设 \((f_j)_{j\in \mathbf{J}}\) 为 F 的标准正交基。令 \(\lambda_{j} = \| u^{*}(f_{j})\|^{2}\)（对所有 \(j\in \mathbf{J}\)）。由于 \(u\) 属于 \(\mathcal{L}^2 (\mathrm{E};\mathrm{F})\)，我们有 \(u^{*}\in \mathcal{L}^{2}(\mathrm{F};\mathrm{E})\)，因而 \(\sum_{j}\lambda_{j} < + \infty\)。此外，我们有

\[
Q (x) = \left\| u (x) \right\| ^ {2} = \sum_ {j} \left| \langle u ^ {*} (f _ {j}) | x \rangle \right| ^ {2}\tag{43}
\]

由 Parseval 公式（V, p.~22）及伴随映射的定义（V, p.~38）。对每个 \(x \in \mathbf{B}\)，由 Cauchy-Schwarz 不等式有 \(|\langle u^{*}(f_{j})|x\rangle|^{2} \leqslant \lambda_{j}\)；因此，公式 (43) 中和的收敛在 B 上是一致的，从而得到引理 4（GT, X, § 1, No.~6）。

引理 5. --- 设 \(\mathbf{E}_1\) 为 \(\mathbf{E}\) 的闭向量子空间，在 \(u^*u\) 下稳定。若 \(\mathbf{E}_1 \neq \{0\}\)，则在 \(\mathbf{E}_1 \cap \mathbf{P}\) 中存在一个范数为 1 的向量。

由于 B 是弱紧的，B 的弱闭子空间 \(B \cap E_{1}\) 也是弱紧的。因而存在（GT, IV, §6, No.~1, 定理 1）一个点 \(x_{0}\)（属于 \(B \cap E_{1}\)），使得 \(Q(x_{0} \geqslant Q(x))\) 对所有 \(x \in B \cap E_{1}\) 成立。若 \(Q(x_{0}) = 0\)，则 \(Q(x) = 0\) 且 \(u(x) = 0\)（对所有 \(x \in B \cap E_{1}\)）。于是 \(E_{1} \subset P\)，引理 5 在这一情形成立。

现在设 \(\mathbf{Q}(x_0) > 0\)，则 \(x_0 \neq 0\)。由于向量 \(\| x_0 \|^{-1} \cdot x_0\) 属于 \(\mathbf{B} \cap \mathbf{E}_1\)，我们有

\[
\mathrm{Q} \left(x _ {0}\right) \geqslant \mathrm{Q} \left(\left\| x _ {0} \right\| ^ {- 1}. x _ {0}\right) = \mathrm{Q} \left(x _ {0}\right) / \left\| x _ {0} \right\| ^ {2}
\]

即 \(\|x_{0}\|=1\)。我们将证明 \(x_{0}\) 属于 P；设 \(y\in E\) 与 \(x_{0}\) 正交。只需证明 \(u(y)\) 与 \(u(x_{0})\) 正交。但由于 y 是 \(E_{1}\) 中一个向量与一个与 \(E_{1}\) 正交的向量之和，且两者都与 \(x_{0}\) 正交（因为 \(x_{0}\in E_{1}\)），所以只需考虑下面两种情形：

\begin{enumerate}
\def\labelenumi{\alph{enumi})}
\item
  \(y\) 与 \(\mathbf{E}_1\) 正交：由于 \(\mathbf{E}_1\) 在 \(u^* u, u^* u(x_0) \in \mathbf{E}_1\) 下稳定，因而 \(0 = \langle y | u^* u(x_0) \rangle = \langle u(y) | u(x_0) \rangle\)。
\item
  \(y\) 属于 \(\mathbf{E}_1\)：对所有 \(t \in \mathbf{R}\)，向量 \(x(t) = (x_0 + ty) / \| x_0 + ty \|\) 属于 \(\mathbf{B} \cap \mathbf{E}_1\)。我们有 \(Q(xt) = f(t) / g(t)\)，其中
\end{enumerate}

\[
f (t) = \left\| u (x _ {0}) \right\| ^ {2} + 2 t \mathscr {R} \left\langle u (x _ {0}) | u (y) \right\rangle + t ^ {2} \left\| u (y) \right\| ^ {2}
\]

\[
g (t) = 1 + t ^ {2} \| y \| ^ {2}.
\]

根据 \(x_0\) 的定义，我们有 \(\mathrm{Q}(x(0)) \geqslant \mathrm{Q}(x(t))\)（对所有实数 \(t\)），因而 \(\frac{d}{dt} \mathrm{Q}(x(t))\) 在 \(t = 0\) 处为零。但 \(f(0) = \| u(x_0)\|^2\)，\(g(0) = 1\)，\(f'(0) = 2\mathcal{R}\langle u(x_0)|u(y)\rangle\)，\(g'(0) = 0\)。由于

\[
\frac {d}{d t} \mathrm{Q} (x (t)) = \frac {f ^ {\prime} (t) g (t) - f (t) g ^ {\prime} (t)}{g (t) ^ {2}},
\]

我们断定 \(f'(0)=0\)，即 \(\mathcal{R}\langle u(x_{0})|u(y)\rangle=0\)。当 K=R 时，\(u(x_{0})\) 与 \(u(y)\) 正交；当 K=C 时，向量 iy 属于 \(E_{1}\) 且与 \(x_{0}\) 正交，因而 \(\mathcal{I}\langle u(x_{0})|u(y)\rangle=-\mathcal{R}\langle u(x_{0})|u(iy)\rangle=0\)，最终 \(u(x_{0})\) 与 \(u(y)\) 正交。这就证明了引理 5。

现在我们证明定理 2。应用 S, III, §4, No.~5 的定理 1，如同 V, p.~23 那样，我们看到存在一个集 S，它在 E 的包含于 P 的标准正交子集中是极大的。令 \(E_{1}\) 为与 S 正交的所有向量组成的集。设 \(y \in E_{1}\)；若 \(x \in S\)，则向量 x 与 y 正交，且由于 \(S \subset P\)，我们断定 \(u(x)\) 与 \(u(y)\) 正交；于是

\[
\langle x | u ^ {*} u (y) \rangle = \langle u (x) | u (y) \rangle = 0
\]

并且 \(u^{*}u(y)\) 与 S 正交。因而 \(E_{1}\) 在 \(u^{*}u\) 下稳定。倘若 \(E_{1} \neq \{0\}\)，就会存在一个范数为 1 的向量 \(x\)（属于 \(E_{1} \cap P\)）（引理 5），而 \(S \cup \{x\}\) 将是 E 的包含于 P 的标准正交子集。这与 S 的极大性矛盾。因而 \(E_{1} = \{0\}\)，S 是 E 的标准正交基。证毕。

推论 1. --- 设 v 为希尔伯特空间 E 的具有有限迹的连续正自同态。存在 E 的标准正交基 \((e_{i})_{i\in I}\) 及正实数的可和族 \((\lambda_{i})_{i\in I}\)，使得 \(v(e_{i}) = \lambda_{i}e_{i}\) 对所有 \(i \in I\) 成立。

令 \(\Phi(x, y) = \langle x | v(y) \rangle\)（对 E 中的 \(x, y\)）。则 \(\Phi\) 是 E 上的正埃尔米特型。存在（V, p.~8, 推论）一个希尔伯特空间 F 及一个从 E 到 F 中的连续线性映射 \(u\)，使得 \(\Phi(x, y) = \langle u(x) | u(y) \rangle\)（对 E 中的 \(x, y\)）。换言之，我们有 \(v = u^{*}u\)。根据定义 9（V, p.~52），\(u\) 是从 E 到 F 中的 Hilbert-Schmidt 映射。由定理 2，存在 E 的标准正交基 \((e_{i})_{i \in I}\)，使得向量 \(u(e_{i})\) 两两正交。设 \(i \in I\)；对 I 中每个 \(j \neq i\)，我们有

\[
\langle e _ {j} | v (e _ {i}) \rangle = \langle u (e _ {j}) | u (e _ {i}) \rangle = 0
\]

因而 \(v(e_{i})\) 与 \(e_{i}\) 成比例，且具有 \(\lambda_{i}e_{i}\) 的形式，其中 \(\lambda_{i}=\langle e_{i}|v(e_{i})\rangle\)；于是

\[
\lambda_ {i} \geqslant 0 \quad \text { and } \quad \sum_ {i \in I} \lambda_ {i} = \operatorname{Tr} (v) <   + \infty .
\]

推论 2. --- 设 E 为希尔伯特空间。则 \(\mathcal{L}^{1}(\mathrm{E})\subset \mathcal{L}^{2}(\mathrm{E})\)

实的情形可通过标量扩张归结为复的情形；因此我们可以设 \(\mathbf{K} = \mathbf{C}\)。

由于 \(\mathcal{L}^{2}(E)\) 是 \(\mathcal{L}(E)\) 的向量子空间，只需证明 E 的每个具有有限迹的连续正自同态 v 属于 \(\mathcal{L}^{2}(E)\)。用推论 1 的记号，我们有

\[
\sum_ {i \in \mathbf {I}} \left\| v (e _ {i}) \right\| ^ {2} = \sum_ {i \in \mathbf {I}} \lambda_ {i} ^ {2} \leqslant (\sum_ {i} \lambda_ {i}) ^ {2} <   + \infty .
\]

推论 3. --- 设 v 为希尔伯特空间 E 的具有有限迹的连续正自同态。存在 E 的一个正 Hilbert-Schmidt 自同态 w，使得 \(v = w^{2}\) 且 v 与 w 交换。

用推论 1 的记号，只需考虑自同态 \(w\)，它把向量 \(\sum_{i\in I}\xi_i e_i\) 变为向量 \(\sum_{i}\lambda_i^{1 / 2}\xi_i e_i\)。

注. --- 用定理 2 的记号，令 J 为所有 \(i \in \mathbf{I}\)（满足 \(u(e_i) \neq 0\)）组成的集。对所有 \(i \in \mathbf{J}\)，令 \(\lambda_i = \| u(e_i)\|\) 及 \(f_i = \lambda_i^{-1}u(e_i)\)。则 \((e_i)_{i \in \mathbf{J}}(\text{resp. } (f_i)_{i \in \mathbf{J})}\) 是 \(u\) 的初始（相应地终）子空间的标准正交基，我们有 \(u(e_i) = \lambda_i f_i\)（对所有 \(i \in \mathbf{J}\)），且 \(\sum_{i \in \mathbf{J}} \lambda_i^2 = \| u\|_2^2\) 是有限的。

